% {Describir detalladamente el problema, sin incluir la solución.
%  Incluir referencias a todos los datos manejados.}
% [Mínimo media cuartilla, máximo dos]
La toma de decisiones en un equipo durante el desarrollo de un partido en la WNBA es un factor que determina
el éxito o fracaso del juego. El problema radica en que las decisiones de sustitución se fundamentan en juicios
subjetivos, como en la intuición del entrenador, experiencias previas o percepciones en un momento decisivo
del partido, en lugar de basarse en modelos cuantitativos dependientes del contexto [1]$\leftarrow$ CHECAR REFERENCIA.
\\
En el baloncesto profesional el análisis de datos como los "play-by-play", ha evolucionado constantemente;
entre ellos uno de los estudios recientes, aplican técnicas de aprendizaje automático para la predicción de
resultados en partidos de baloncesto; sin embargo, muy pocos estudios se centran en el estudio de la sustitución
de jugadores a lo largo del partido de basquetbol. En este contexto el artículo de Bonato y Hinds (2026) da un
punto de inflexión, en dicho estudio se centran en equipos femeninos universitarios de la primera división,
haciendo uso de grafos dirigidos para el modelado de las competiciones pases y acciones defensivas como
redes, haciendo uso del Algoritmo PageRank, que originalmente era usado para ordenar páginas web, como
también el uso del algoritmo node2vec donde se hacen representaciones de grafos para poder hacer uso de
modelos de machine learning tradicionales \cite{bonato2026networkanalysislinkprediction}.